\section{Unüberwachtes Lernen}
\setcounter{aufgabe}{0}

% ===================================================================
% Aufgaben zu Unueberwachtem Lernen
% Themen: MDS / Stress-Funktion, SMACOF / Guttman-Transformation,
%         t-SNE (Aehnlichkeiten, Symmetrisierung, KL-Divergenz),
%         k-Means (WCSS, Zuordnung, Update, Elbow)
% ===================================================================

% -------------------------------------------------------------------
% 1) Distanzmatrix und einfache Projektion
% -------------------------------------------------------------------
\aufgabe{
Gegeben seien vier Punkte im $\mathbb{R}^2$:
\[
  \mathbf{p}_1 = \begin{pmatrix} 0 \\ 0 \end{pmatrix}, \quad
  \mathbf{p}_2 = \begin{pmatrix} 3 \\ 0 \end{pmatrix}, \quad
  \mathbf{p}_3 = \begin{pmatrix} 0 \\ 4 \end{pmatrix}, \quad
  \mathbf{p}_4 = \begin{pmatrix} 3 \\ 4 \end{pmatrix}.
\]

\begin{enumerate}[label=\alph*)]
  \item Berechnen Sie die euklidische Distanzmatrix
        $D = (d_{ij})$ mit $d_{ij} = \lVert \mathbf{p}_i - \mathbf{p}_j \rVert_2$.
  \item Projizieren Sie die Punkte durch Weglassen der zweiten Koordinate
        auf die $x_1$-Achse. Welche Distanzen werden dabei verfälscht?
  \item Zeigen Sie an einem Paar, dass diese Projektion \textbf{nicht}
        abstandserhaltend ist.
\end{enumerate}
}

\loesung{
\begin{enumerate}[label=\alph*)]
  \item Die Distanzmatrix ist symmetrisch mit Nulldiagonale:
  \[
    D =
    \begin{pmatrix}
      0 & 3 & 4 & 5\\
      3 & 0 & 5 & 4\\
      4 & 5 & 0 & 3\\
      5 & 4 & 3 & 0
    \end{pmatrix}
  \]
  Beispielrechnung:
  $d_{14} = \sqrt{3^2 + 4^2} = 5$, $d_{23} = \sqrt{3^2+4^2} = 5$.

  \item Nach der Projektion gilt $\tilde{p}_i = (p_{i,1})$, also
  \[
    \tilde{p}_1 = 0, \quad \tilde{p}_2 = 3, \quad
    \tilde{p}_3 = 0, \quad \tilde{p}_4 = 3.
  \]
  Punkte, die sich nur in der zweiten Koordinate unterscheiden, fallen
  zusammen: $\mathbf{p}_1, \mathbf{p}_3$ auf $0$ und
  $\mathbf{p}_2, \mathbf{p}_4$ auf $3$.

  \item Für das Paar $(1,3)$ gilt im Original $d_{13} = 4$, in der Projektion
  aber
  \[
    \tilde{d}_{13} = |0 - 0| = 0 \neq 4.
  \]
  Zwei verschiedene Punkte werden also auf denselben Bildpunkt abgebildet.
  Die Projektion ist damit nicht abstandserhaltend.
\end{enumerate}
}

% -------------------------------------------------------------------
% 2) MDS: Stress-Funktion
% -------------------------------------------------------------------
\aufgabe{
Für drei Objekte seien die Zielabstände (Unähnlichkeiten)
\[
  \delta_{12} = 2, \qquad \delta_{13} = 4, \qquad \delta_{23} = 3
\]
gegeben. Alle Gewichte seien $w_{ij} = 1$. Eine Konfiguration $X$ auf der
Geraden liefert die Bildpunkte
\[
  x_1 = 0, \qquad x_2 = 1, \qquad x_3 = 5 .
\]
Der Stress ist
\[
  \sigma(X) = \sum_{i<j} w_{ij}\bigl(\delta_{ij} - d_{ij}(X)\bigr)^2 .
\]

\begin{enumerate}[label=\alph*)]
  \item Berechnen Sie die Abstände $d_{ij}(X) = |x_i - x_j|$.
  \item Berechnen Sie den Stress $\sigma(X)$.
  \item Verschieben Sie $x_2$ auf $x_2 = 2$. Sinkt der Stress?
\end{enumerate}
}

\loesung{
\begin{enumerate}[label=\alph*)]
  \item
  \[
    d_{12} = |0-1| = 1, \qquad
    d_{13} = |0-5| = 5, \qquad
    d_{23} = |1-5| = 4 .
  \]

  \item
  \[
    \sigma(X) = (2-1)^2 + (4-5)^2 + (3-4)^2 = 1 + 1 + 1 = 3 .
  \]

  \item Mit $x_2 = 2$:
  \[
    d_{12} = 2, \qquad d_{23} = 3, \qquad d_{13} = 5 \text{ (unverändert)} .
  \]
  \[
    \sigma(X') = (2-2)^2 + (4-5)^2 + (3-3)^2 = 0 + 1 + 0 = 1 < 3 .
  \]
  Der Stress sinkt von $3$ auf $1$, die Konfiguration ist also besser.
\end{enumerate}
}

% -------------------------------------------------------------------
% 3) SMACOF: Guttman-Transformation
% -------------------------------------------------------------------
\aufgabe{
Für $n = 3$ Punkte auf der Geraden seien die Zielabstände
\[
  \delta_{12} = 1, \qquad \delta_{13} = 3, \qquad \delta_{23} = 2
\]
gegeben ($w_{ij} = 1$). Die aktuelle Konfiguration ist
\[
  Z = (z_1, z_2, z_3)^{\mathsf{T}} = (0,\; 2,\; 3)^{\mathsf{T}} .
\]
Ein SMACOF-Schritt lautet
$X^{(k+1)} = \tfrac{1}{n}\,B(Z)\,Z$ mit
\[
  B(Z)_{ij} =
  \begin{cases}
    -\dfrac{\delta_{ij}}{d_{ij}(Z)}, & i \neq j,\ d_{ij}(Z) \neq 0,\\[0.3cm]
    -\displaystyle\sum_{l \neq i} B(Z)_{il}, & i = j.
  \end{cases}
\]

\begin{enumerate}[label=\alph*)]
  \item Berechnen Sie die Abstände $d_{ij}(Z) = |z_i - z_j|$.
  \item Stellen Sie die Matrix $B(Z)$ auf.
  \item Berechnen Sie die neue Konfiguration
        $X^{(1)} = \tfrac{1}{3}\,B(Z)\,Z$.
\end{enumerate}
}

\loesung{
\begin{enumerate}[label=\alph*)]
  \item
  \[
    d_{12} = 2, \qquad d_{13} = 3, \qquad d_{23} = 1 .
  \]

  \item Außerdiagonale Einträge $-\delta_{ij}/d_{ij}$:
  \[
    B_{12} = -\tfrac{1}{2}, \qquad
    B_{13} = -\tfrac{3}{3} = -1, \qquad
    B_{23} = -\tfrac{2}{1} = -2 .
  \]
  Diagonale als negative Zeilensumme der Außendiagonalen:
  \[
    B_{11} = \tfrac{1}{2} + 1 = \tfrac{3}{2}, \quad
    B_{22} = \tfrac{1}{2} + 2 = \tfrac{5}{2}, \quad
    B_{33} = 1 + 2 = 3 .
  \]
  \[
    B(Z) =
    \begin{pmatrix}
      \tfrac{3}{2} & -\tfrac{1}{2} & -1\\[0.15cm]
      -\tfrac{1}{2} & \tfrac{5}{2} & -2\\[0.15cm]
      -1 & -2 & 3
    \end{pmatrix}
  \]

  \item Zunächst $B(Z)\,Z$ mit $Z = (0,2,3)^{\mathsf{T}}$:
  \[
    B(Z)\,Z =
    \begin{pmatrix}
      \tfrac{3}{2}\cdot 0 - \tfrac{1}{2}\cdot 2 - 1\cdot 3\\[0.1cm]
      -\tfrac{1}{2}\cdot 0 + \tfrac{5}{2}\cdot 2 - 2\cdot 3\\[0.1cm]
      -1\cdot 0 - 2\cdot 2 + 3\cdot 3
    \end{pmatrix}
    =
    \begin{pmatrix} -4 \\ -1 \\ 5 \end{pmatrix}
  \]
  \[
    X^{(1)} = \tfrac{1}{3}
    \begin{pmatrix} -4 \\ -1 \\ 5 \end{pmatrix}
    =
    \begin{pmatrix} -1{,}\overline{3} \\ -0{,}\overline{3} \\ 1{,}\overline{6} \end{pmatrix}
  \]
  Die neuen Abstände sind $|x_1-x_2| = 1$, $|x_2-x_3| = 2$,
  $|x_1-x_3| = 3$, also \textbf{exakt} die Zielabstände: Der Stress ist
  nach einem Schritt bereits $0$.
\end{enumerate}
}

% -------------------------------------------------------------------
% 4) t-SNE: bedingte Aehnlichkeiten im Ausgangsraum
% -------------------------------------------------------------------
\aufgabe{
Im Ausgangsraum wird die bedingte Ähnlichkeit über eine Normalverteilung
um den Punkt $\mathbf{x}_i$ definiert:
\[
  p_{j|i} =
  \frac{\exp\!\bigl(-\lVert \mathbf{x}_i - \mathbf{x}_j \rVert^2 / 2\sigma_i^2\bigr)}
       {\sum_{k \neq i} \exp\!\bigl(-\lVert \mathbf{x}_i - \mathbf{x}_k \rVert^2 / 2\sigma_i^2\bigr)} .
\]
Betrachtet werde $i = 1$ mit den quadrierten Abständen
\[
  \lVert \mathbf{x}_1 - \mathbf{x}_2 \rVert^2 = 2, \qquad
  \lVert \mathbf{x}_1 - \mathbf{x}_3 \rVert^2 = 8, \qquad
  \lVert \mathbf{x}_1 - \mathbf{x}_4 \rVert^2 = 8 ,
\]
und der Bandbreite $\sigma_1 = 1$.

\begin{enumerate}[label=\alph*)]
  \item Berechnen Sie die unskalierten Ähnlichkeiten
        $\exp(-\lVert \mathbf{x}_1 - \mathbf{x}_j \rVert^2 / 2)$ für
        $j = 2,3,4$.
  \item Berechnen Sie $p_{2|1}$, $p_{3|1}$ und $p_{4|1}$.
  \item Prüfen Sie, dass $\sum_{j \neq 1} p_{j|1} = 1$ gilt, und
        interpretieren Sie das Ergebnis.
\end{enumerate}
Hinweis: $e^{-1} \approx 0{,}368$, $e^{-4} \approx 0{,}0183$.
}

\loesung{
\begin{enumerate}[label=\alph*)]
  \item Mit $2\sigma_1^2 = 2$:
  \[
    s_2 = e^{-2/2} = e^{-1} \approx 0{,}368, \qquad
    s_3 = s_4 = e^{-8/2} = e^{-4} \approx 0{,}0183 .
  \]

  \item Normierung über die Summe
  $s_2 + s_3 + s_4 \approx 0{,}368 + 0{,}0183 + 0{,}0183 = 0{,}4046$:
  \[
    p_{2|1} \approx \frac{0{,}368}{0{,}4046} \approx 0{,}910,
  \]
  \[
    p_{3|1} = p_{4|1} \approx \frac{0{,}0183}{0{,}4046} \approx 0{,}045 .
  \]

  \item Summe:
  \[
    p_{2|1} + p_{3|1} + p_{4|1} \approx 0{,}910 + 0{,}045 + 0{,}045 = 1 .
  \]
  Der nahe Nachbar $\mathbf{x}_2$ erhält fast die gesamte
  Wahrscheinlichkeitsmasse, die fernen Punkte nur einen kleinen Anteil.
  t-SNE gewichtet damit die \textbf{lokale} Nachbarschaft stark.
\end{enumerate}
}

% -------------------------------------------------------------------
% 5) t-SNE: Zielraum-Aehnlichkeiten und KL-Divergenz
% -------------------------------------------------------------------
\aufgabe{
Im Zielraum verwendet t-SNE eine Student-t-Verteilung mit einem
Freiheitsgrad:
\[
  q_{ij} =
  \frac{\bigl(1 + \lVert \mathbf{y}_i - \mathbf{y}_j \rVert^2\bigr)^{-1}}
       {\sum_{k \neq l}\bigl(1 + \lVert \mathbf{y}_k - \mathbf{y}_l \rVert^2\bigr)^{-1}} .
\]
Drei Punkte im Zielraum haben die quadrierten Abstände
\[
  \lVert \mathbf{y}_1 - \mathbf{y}_2 \rVert^2 = 0, \qquad
  \lVert \mathbf{y}_1 - \mathbf{y}_3 \rVert^2 = 1, \qquad
  \lVert \mathbf{y}_2 - \mathbf{y}_3 \rVert^2 = 3 .
\]

\begin{enumerate}[label=\alph*)]
  \item Berechnen Sie die unnormierten Werte
        $(1 + \lVert \mathbf{y}_i - \mathbf{y}_j \rVert^2)^{-1}$ und daraus
        $q_{12}$, $q_{13}$, $q_{23}$.
  \item Gegeben sei die Zielverteilung
        $p_{12} = 0{,}6$, $p_{13} = 0{,}3$, $p_{23} = 0{,}1$. Berechnen Sie
        die Kostenfunktion
        \[
          C = \sum_{i < j} p_{ij}\,\log \frac{p_{ij}}{q_{ij}} .
        \]
        Interpretieren Sie den größten Summanden.
\end{enumerate}
Hinweis: natürlicher Logarithmus, $\log(1{,}05) \approx 0{,}050$,
$\log(0{,}70) \approx -0{,}357$.
}

\loesung{
\begin{enumerate}[label=\alph*)]
  \item Unnormierte Werte:
  \[
    t_{12} = \frac{1}{1+0} = 1, \qquad
    t_{13} = \frac{1}{1+1} = 0{,}5, \qquad
    t_{23} = \frac{1}{1+3} = 0{,}25 .
  \]
  Summe über alle Paare (jedes Paar einfach): $1 + 0{,}5 + 0{,}25 = 1{,}75$.
  \[
    q_{12} = \frac{1}{1{,}75} \approx 0{,}571, \quad
    q_{13} = \frac{0{,}5}{1{,}75} \approx 0{,}286, \quad
    q_{23} = \frac{0{,}25}{1{,}75} \approx 0{,}143 .
  \]

  \item Einsetzen:
  \[
    C = 0{,}6\log\frac{0{,}6}{0{,}571}
      + 0{,}3\log\frac{0{,}3}{0{,}286}
      + 0{,}1\log\frac{0{,}1}{0{,}143} .
  \]
  \[
    C \approx 0{,}6\cdot 0{,}050 + 0{,}3\cdot 0{,}048
      + 0{,}1\cdot(-0{,}357) \approx 0{,}030 + 0{,}014 - 0{,}036 = 0{,}008 .
  \]
  Der kleine Gesamtwert zeigt, dass $Q$ die Verteilung $P$ bereits gut
  nachbildet. Der Term für das Paar $(1,2)$ trägt am stärksten positiv bei,
  da dort $p_{12}$ am größten ist und Abweichungen bei nahen Paaren
  besonders gewichtet werden.
\end{enumerate}
}

% -------------------------------------------------------------------
% 6) k-Means: Zuordnung und Update (ein Schritt)
% -------------------------------------------------------------------
\aufgabe{
Gegeben seien die eindimensionalen Datenpunkte
\[
  x_1 = 1, \quad x_2 = 2, \quad x_3 = 4, \quad
  x_4 = 7, \quad x_5 = 8
\]
sowie die Startzentren $\mu_1 = 1$ und $\mu_2 = 5$ ($K = 2$).

\begin{enumerate}[label=\alph*)]
  \item Ordnen Sie jeden Punkt dem nächstgelegenen Zentrum zu.
  \item Aktualisieren Sie die Zentren als Mittelwert der zugeordneten
        Punkte.
  \item Führen Sie einen zweiten Zuordnungsschritt durch. Ändert sich die
        Zuordnung noch?
\end{enumerate}
}

\loesung{
\begin{enumerate}[label=\alph*)]
  \item Abstand zu $\mu_1 = 1$ bzw. $\mu_2 = 5$:
  \[
    \begin{array}{c|c|c|c}
      x & |x-1| & |x-5| & \text{Cluster}\\\hline
      1 & 0 & 4 & C_1\\
      2 & 1 & 3 & C_1\\
      4 & 3 & 1 & C_2\\
      7 & 6 & 2 & C_2\\
      8 & 7 & 3 & C_2
    \end{array}
  \]
  Also $C_1 = \{1, 2\}$ und $C_2 = \{4, 7, 8\}$.

  \item Neue Zentren:
  \[
    \mu_1 = \frac{1+2}{2} = 1{,}5, \qquad
    \mu_2 = \frac{4+7+8}{3} = \frac{19}{3} \approx 6{,}33 .
  \]

  \item Zweite Zuordnung mit $\mu_1 = 1{,}5$, $\mu_2 \approx 6{,}33$:
  \[
    \begin{array}{c|c|c|c}
      x & |x-1{,}5| & |x-6{,}33| & \text{Cluster}\\\hline
      1 & 0{,}5 & 5{,}33 & C_1\\
      2 & 0{,}5 & 4{,}33 & C_1\\
      4 & 2{,}5 & 2{,}33 & C_2\\
      7 & 5{,}5 & 0{,}67 & C_2\\
      8 & 6{,}5 & 1{,}67 & C_2
    \end{array}
  \]
  Die Zuordnung bleibt $C_1 = \{1,2\}$, $C_2 = \{4,7,8\}$ unverändert. Das
  Verfahren ist \textbf{konvergiert}.
\end{enumerate}
}

% -------------------------------------------------------------------
% 7) k-Means: WCSS und Vergleich zweier Zuordnungen
% -------------------------------------------------------------------
\aufgabe{
Für die Punkte $x_1 = 1$, $x_2 = 2$, $x_3 = 4$, $x_4 = 7$, $x_5 = 8$ ist
die Zielfunktion
\[
  J = \sum_{k=1}^{K} \sum_{x \in C_k} (x - \mu_k)^2 .
\]

\begin{enumerate}[label=\alph*)]
  \item Berechnen Sie $J$ für die Zuordnung $C_1 = \{1,2\}$,
        $C_2 = \{4,7,8\}$ mit den zugehörigen Mittelwerten.
  \item Berechnen Sie $J$ für die alternative Zuordnung
        $C_1 = \{1,2,4\}$, $C_2 = \{7,8\}$.
  \item Welche Zuordnung ist im Sinne von k-Means besser?
\end{enumerate}
}

\loesung{
\begin{enumerate}[label=\alph*)]
  \item Mittelwerte $\mu_1 = 1{,}5$, $\mu_2 = \tfrac{19}{3} \approx 6{,}33$:
  \[
    J_1 = \bigl((1-1{,}5)^2 + (2-1{,}5)^2\bigr)
        + \bigl((4-6{,}33)^2 + (7-6{,}33)^2 + (8-6{,}33)^2\bigr)
  \]
  \[
    J_1 = (0{,}25 + 0{,}25) + (5{,}43 + 0{,}45 + 2{,}79)
        \approx 0{,}5 + 8{,}67 = 9{,}17 .
  \]

  \item Mittelwerte $\mu_1 = \tfrac{1+2+4}{3} = \tfrac{7}{3} \approx 2{,}33$,
  $\mu_2 = 7{,}5$:
  \[
    J_2 = \bigl((1-2{,}33)^2 + (2-2{,}33)^2 + (4-2{,}33)^2\bigr)
        + \bigl((7-7{,}5)^2 + (8-7{,}5)^2\bigr)
  \]
  \[
    J_2 = (1{,}78 + 0{,}11 + 2{,}79) + (0{,}25 + 0{,}25)
        \approx 4{,}67 + 0{,}5 = 5{,}17 .
  \]

  \item Wegen $J_2 \approx 5{,}17 < 9{,}17 \approx J_1$ ist die zweite
  Zuordnung im Sinne von k-Means \textbf{besser}: Ihre Cluster sind
  kompakter (kleinere WCSS).
\end{enumerate}
}

% -------------------------------------------------------------------
% 8) k-Means im R^2
% -------------------------------------------------------------------
\aufgabe{
Gegeben seien folgende Punkte in $\mathbb{R}^2$:
\[
  \mathbf{x}_1 = \begin{pmatrix} 1 \\ 1 \end{pmatrix}, \;
  \mathbf{x}_2 = \begin{pmatrix} 1 \\ 2 \end{pmatrix}, \;
  \mathbf{x}_3 = \begin{pmatrix} 5 \\ 4 \end{pmatrix}, \;
  \mathbf{x}_4 = \begin{pmatrix} 6 \\ 5 \end{pmatrix}, \;
  \mathbf{x}_5 = \begin{pmatrix} 6 \\ 3 \end{pmatrix}
\]
mit Startzentren
$\boldsymbol{\mu}_1 = \begin{pmatrix} 1 \\ 1 \end{pmatrix}$ und
$\boldsymbol{\mu}_2 = \begin{pmatrix} 5 \\ 5 \end{pmatrix}$.

\begin{enumerate}[label=\alph*)]
  \item Ordnen Sie die Punkte anhand des \textbf{quadrierten} euklidischen
        Abstands zu (die Wurzel ist für den Vergleich nicht nötig).
  \item Aktualisieren Sie beide Zentren.
  \item Berechnen Sie die resultierende WCSS $J$.
\end{enumerate}
}

\loesung{
\begin{enumerate}[label=\alph*)]
  \item Quadrierte Abstände zu $\boldsymbol{\mu}_1 = (1,1)$ und
  $\boldsymbol{\mu}_2 = (5,5)$:
  \[
    \begin{array}{c|c|c|c}
      \mathbf{x} & \lVert\cdot - \boldsymbol{\mu}_1\rVert^2
                 & \lVert\cdot - \boldsymbol{\mu}_2\rVert^2 & \text{Cluster}\\\hline
      (1,1) & 0 & 32 & C_1\\
      (1,2) & 1 & 25 & C_1\\
      (5,4) & 25 & 1 & C_2\\
      (6,5) & 41 & 1 & C_2\\
      (6,3) & 29 & 5 & C_2
    \end{array}
  \]
  Also $C_1 = \{\mathbf{x}_1, \mathbf{x}_2\}$,
  $C_2 = \{\mathbf{x}_3, \mathbf{x}_4, \mathbf{x}_5\}$.

  \item Neue Zentren als komponentenweiser Mittelwert:
  \[
    \boldsymbol{\mu}_1 = \begin{pmatrix} 1 \\ 1{,}5 \end{pmatrix},
    \qquad
    \boldsymbol{\mu}_2 =
    \begin{pmatrix} \tfrac{5+6+6}{3} \\ \tfrac{4+5+3}{3} \end{pmatrix}
    = \begin{pmatrix} 5{,}\overline{6} \\ 4 \end{pmatrix} .
  \]

  \item WCSS mit den neuen Zentren:
  \[
    \text{Cluster } 1: (0 + 0{,}25) + (0 + 0{,}25) = 0{,}5
  \]
  \[
    \text{Cluster } 2:
    \underbrace{(0{,}\overline{4} + 0)}_{(5,4)}
    + \underbrace{(0{,}\overline{1} + 1)}_{(6,5)}
    + \underbrace{(0{,}\overline{1} + 1)}_{(6,3)}
    \approx 2{,}67
  \]
  \[
    J = 0{,}5 + 2{,}67 \approx 3{,}17 .
  \]
\end{enumerate}
}

% -------------------------------------------------------------------
% 9) k-Means: Elbow-Methode
% -------------------------------------------------------------------
\aufgabe{
Bei der Elbow-Methode trägt man die WCSS $J(k)$ gegen die Clusterzahl $k$
auf und sucht den \textbf{Knick}. Gemessen wurden:
\[
  \begin{array}{c|c|c|c|c|c|c}
    k    & 1 & 2 & 3 & 4 & 5 & 6\\\hline
    J(k) & 480 & 180 & 70 & 55 & 45 & 38
  \end{array}
\]

\begin{enumerate}[label=\alph*)]
  \item Berechnen Sie die absolute Abnahme $\Delta J(k) = J(k-1) - J(k)$
        für $k = 2, \dots, 6$.
  \item Bestimmen Sie den Knick (Elbow) und begründen Sie die Wahl von $k$.
  \item Warum ist es nicht sinnvoll, einfach das kleinste $J$ zu wählen?
\end{enumerate}
}

\loesung{
\begin{enumerate}[label=\alph*)]
  \item
  \[
    \begin{array}{c|c|c|c|c|c}
      k & 2 & 3 & 4 & 5 & 6\\\hline
      \Delta J(k) & 300 & 110 & 15 & 10 & 7
    \end{array}
  \]

  \item Von $k=1$ auf $2$ und von $2$ auf $3$ sinkt $J$ stark
  ($300$ bzw. $110$), danach nur noch geringfügig ($15, 10, 7$). Der Knick
  liegt bei $k = 3$: Zusätzliche Cluster bringen kaum noch Verbesserung.
  Man wählt daher $\mathbf{k = 3}$.

  \item $J$ sinkt mit wachsendem $k$ \textbf{monoton} und wird für $k = n$
  (jeder Punkt ein eigenes Cluster) sogar $0$. Das kleinste $J$ liefert also
  keine sinnvolle Clusterzahl, sondern führt zu Überanpassung. Gesucht ist
  der beste Kompromiss zwischen Kompaktheit und Anzahl der Cluster.
\end{enumerate}
}